\section{Gap calculations for transfers and applications}
\label{app:transfer-gap-calculations}

This appendix gives the parameter choices used in
Section~\ref{subsec:bounded-degree-consequences} and the regularization
and subdivision corollaries, followed by the state-space application.
In each sparsification calculation,
$\tau>0$ is the fixed error parameter in the appropriate general
sparsification theorem, and $r,r'\geq1$ are approximation factors,
with $r<r'$ whenever both occur. All source enlargements are those described
in Section~\ref{subsec:bounded-degree-consequences}; constants are
chosen before the source order is enlarged.
In the NP calculations, $F$ denotes the normalized source system from
Theorem~\ref{thm:xor-source}, with $m\geq1$ equations. In the UGC
layer calculations, $G$ is the independent-set source and $n=|V(G)|$.
In the SSE calculations, let $G^*$ be the host from
Lemma~\ref{lem:sse-unweighted-gap}, applied with source error
$\varepsilon=\eta$ and $q=\lceil2/\eta\rceil$, and let $c>0$ be its
absolute constant; then
$N=|V(\mathcal L(G^*))|=|E(G^*)|$ is the dense output order.

\subsection{Bounded-degree treedepth}
\label{app:transfer-td}
Use the constructions and factor descriptions in
Corollary~\ref{cor:td-bounded-degree}.

\subsubsection{NP-hardness}
\label{app:transfer-td-np}
Deleting $z_A,z_B,z_C$ lowers treedepth by at most $3$, by
Lemma~\ref{lem:restore-vertices}. The resulting graph has order
$18m$, so sparsification lowers treedepth by at most $18\tau m$.
The resulting YES bound is $(17/2+\epsilon)m+1$, and the NO lower
bound is $(9-\epsilon)m-2-18\tau m$. Fix a target factor
$r<18/17$ and let
$\gamma=9-\epsilon-r(17/2+\epsilon)>0$, using the choice of
$\epsilon$ in Section~\ref{subsec:np-td-gap}. Choose
\begin{equation}\label{eq:bounded-td-np-choices}
 \tau<\gamma/36,\qquad m>2(2+r)/\gamma,
\end{equation}
taking enough source copies for the second condition. Then
\[
(9-\epsilon)m-2-18\tau m
 -r\bigl((17/2+\epsilon)m+1\bigr)
 \overset{\eqref{eq:bounded-td-np-choices}}{\ge}
 \frac{\gamma m}{2}-(2+r)
 \overset{\eqref{eq:bounded-td-np-choices}}{>}0,
\]
so the gap survives both losses.

\subsubsection{UGC hardness}
\label{app:transfer-td-ugc}
Fix $r<r'<4/3$ and the parameters $q,\eta$ chosen in
Section~\ref{subsec:ugc-td-gap} for $r'$. Write $Y$ for
the dense YES upper bound. A layer clique gives
\begin{equation}\label{eq:bounded-td-ugc-layer}
 n\le\td(S_q(G))\le Y\qquad\text{in the YES case}.
\end{equation}
The sparsification loss is therefore at most
$\tau qn\overset{\eqref{eq:bounded-td-ugc-layer}}{\le}\tau qY$.
Choose $\tau<(r'-r)/(q(r+1))$. Lemma~\ref{lem:gap-preservation}
then transfers the dense factor-$r'$ gap to factor $r$.

\subsubsection{SSE hardness}
\label{app:transfer-td-sse}
The dense YES and NO treedepth bounds are $4\eta N$ and $\beta N$,
where $\beta=(c/2)\sqrt\eta$, and sparsification loses at most $\tau N$.
For a target factor $r$, choose $\eta$
small enough that $\gamma:=\beta-4r\eta>0$, which is possible
since $\beta=(c/2)\sqrt{\eta}$. Choose
\begin{equation}\label{eq:bounded-td-sse-choice}
 \tau<\gamma.
\end{equation}
The sparse NO bound exceeds $r$ times the YES bound by at least
${(\gamma-\tau)N\overset{\eqref{eq:bounded-td-sse-choice}}{>}0}$.

\subsection{Bounded-degree treewidth}
\label{app:transfer-tw}
Use the constructions and factor descriptions in
Corollary~\ref{cor:tw-bounded-degree}.

\subsubsection{NP-hardness}
\label{app:transfer-tw-np}
Fix a target factor $r<12/11$. The dense gap calculation in
Section~\ref{subsec:np-treewidth} chooses $\epsilon,R$ so that
\begin{equation}\label{eq:bounded-tw-np-gap}
 \gamma=R(6-\epsilon)-r\bigl(R(11/2+\epsilon)+3\bigr)>0.
\end{equation}
The dense YES and NO bounds are respectively
$R(11/2+\epsilon)m+3m-1$ and $R(6-\epsilon)m-1$,
using Lemma~\ref{lem:tw-replication}.
Since the output order is $9Rm$, sparsification preserves the YES
upper bound and lowers the NO bound by at most $9\tau Rm$. Choose
\begin{equation}\label{eq:bounded-tw-np-choice}
 \tau<\gamma/(18R).
\end{equation}
The sparse NO lower bound minus $r$ times the
sparse YES upper bound is
\[
(\gamma-9\tau R)m+(r-1)
 \overset{\eqref{eq:bounded-tw-np-choice}}{>}
 \frac{\gamma m}{2}+r-1
 \overset{\eqref{eq:bounded-tw-np-gap}}{>}0.
\]

\subsubsection{UGC hardness}
\label{app:transfer-tw-ugc}
Fix $r<r'<4/3$ and the parameters $\zeta,R$ used for the dense
factor-$r'$ gap in Section~\ref{subsec:ugc-tw-gap}.
Let $n$ be the order of the independent-set source $G$, so $P_3(G)$
has $3n$ vertices.
The dense graph has global cliques of size $Rn$, so its treewidth is
at least $Rn-1$. Write $Y$ for the dense YES upper bound, so
\begin{equation}\label{eq:bounded-tw-ugc-scale}
 Y\ge Rn-1.
\end{equation}
Since the graph has order $3Rn$, the sparsification
loss, when $Rn\ge2$, is at most
\[
 3\tau Rn\le6\tau(Rn-1)
 \overset{\eqref{eq:bounded-tw-ugc-scale}}{\le}6\tau Y.
\]
Choose
$\tau<(r'-r)/(6(r+1))$. Lemma~\ref{lem:gap-preservation} then
transfers the dense factor-$r'$ gap to factor $r$.

\subsubsection{SSE hardness}
\label{app:transfer-tw-sse}
For target factor $r$, the sparse NO bound is $\beta N-1-\tau N$
and the YES bound is $4\eta N$, with $\beta=(c/2)\sqrt\eta$.
Fix $\eta$ so that $\gamma:=\beta-4r\eta>0$, which is possible
because $\beta=(c/2)\sqrt\eta$. Choose
\begin{equation}\label{eq:bounded-tw-sse-choices}
 \tau<\gamma/2,\qquad N>2/\gamma,
\end{equation}
taking sufficiently many parallel copies for the second condition. Then
\[
 \beta N-1-\tau N
 \overset{\eqref{eq:bounded-tw-sse-choices}}{>}
 (\beta-\gamma)N=4r\eta N,
\]
so the sparse NO width exceeds $r$ times the YES bound.

\subsection{Bounded-degree total depth}
\label{app:transfer-total-depth}
Use the constructions and factor descriptions in
Corollary~\ref{cor:total-depth-bounded-degree}.

\subsubsection{NP-hardness}
\label{app:transfer-total-np}
Fix $1\le r<r'<63/62$ and the dense graph
$H=H^\Sigma_M(G_F^\Sigma)$ from
Theorem~\ref{thm:np-total-depth-dense-gap}; $F$ has $m$ equations
and $M=\Theta(m)$.
Write $N=99m+3M$ for the output order and $Y$ for the dense YES
upper bound. Proposition~\ref{prop:np-total-depth-parity-identity}
and~\eqref{eq:np-total-depth-added-vertices-bounds} give
\begin{equation}\label{eq:bounded-total-depth-np-dense-lower}
 \totd(H)
 \overset{\eqref{eq:np-total-depth-added-vertices-bounds},\,
             \eqref{eq:np-total-depth-parity-W}}{\ge}
 M(186m+3)\ge Mm.
\end{equation}
Since $M=\Theta(m)$ and $N=\Theta(m)$, fix $c_*>0$ such that
\begin{equation}\label{eq:bounded-total-depth-np-scale}
 Mm\ge c_*N^2,\qquad Y\ge Mm.
\end{equation}
The second inequality follows because $Y$ is an upper bound on
the dense YES optimum. Choose
\begin{equation}\label{eq:bounded-total-depth-np-choice}
 0<\tau<c_*\min\{r'-r,1/2\}.
\end{equation}
Sparsification preserves the YES
upper bound $Y$ and reduces the NO lower bound by at most
\[
\tau N^2
 \overset{\eqref{eq:bounded-total-depth-np-scale}}{\le}
 (\tau/c_*)Mm
 \overset{\eqref{eq:bounded-total-depth-np-choice}}{<}
 (r'-r)Mm
 \overset{\eqref{eq:bounded-total-depth-np-scale}}{\le}(r'-r)Y.
\]
Hence it preserves the factor-$r$ gap.
Moreover, writing $H'$ for the sparse graph gives
\begin{equation}\label{eq:bounded-total-depth-np-sparse}
\totd(H')
 \overset{\text{Thm.~\ref{thm:total-depth-sparsification};\,
                   \eqref{eq:bounded-total-depth-np-dense-lower}}}{\ge}
 Mm-\tau N^2
  \overset{\eqref{eq:bounded-total-depth-np-scale},\,
              \eqref{eq:bounded-total-depth-np-choice}}{\ge}Mm/2.
\end{equation}
This is $\Omega(N^2)$ by~\eqref{eq:bounded-total-depth-np-scale},
as needed for the regular-graph transfer below.

\subsubsection{UGC hardness}
\label{app:transfer-total-ugc}
Fix $r<r'<8/7$ and the dense parameters $q,\eta$ from
Section~\ref{subsec:ugc-total-depth-gap}, with $N=qn$.
Each layer is a clique, so its $n$ vertices lie on one ancestor
path and have depths at least $1,2,\ldots,n$. Summing over all $q$
layers gives the lower bound
\begin{equation}\label{eq:bounded-total-depth-ugc-scale}
L:=q\,\frac{n(n+1)}2\le\totd(S_q(G)),
 \qquad L\ge\frac{qn^2}{2}=\frac{N^2}{2q},\quad N=qn.
\end{equation}
Write $Y$ for the dense YES upper bound. In the YES case, the following
bounds control the sparsification loss:
\begin{equation}\label{eq:bounded-total-depth-ugc-yes}
L\overset{\eqref{eq:bounded-total-depth-ugc-scale}}{\le}
 \totd(S_q(G))\le Y,\qquad \tau N^2
 \overset{\eqref{eq:bounded-total-depth-ugc-scale}}{\le}2q\tau L
 \overset{\eqref{eq:bounded-total-depth-ugc-yes}}{\le}2q\tau Y.
\end{equation}
Choose
$\tau<(r'-r)/(2q(r+1))$. Lemma~\ref{lem:gap-preservation} then
transfers the dense factor-$r'$ gap to factor $r$.

\subsubsection{SSE hardness}
\label{app:transfer-total-sse}
The dense YES and NO total-depth bounds are $4\eta N^2$ and
$\beta N^2/2$, where $\beta=(c/2)\sqrt\eta$. Sparsification loses
at most $\tau N^2$.
For a target factor $r$, choose $\eta$
small enough that $\gamma:=\beta/2-4r\eta>0$, which is possible
since $\beta=(c/2)\sqrt{\eta}$. Choose
\begin{equation}\label{eq:bounded-total-depth-sse-choice}
 \tau<\gamma.
\end{equation}
After
sparsification, the NO lower bound minus $r$ times the YES upper
bound is
$(\gamma-\tau)N^2\overset{\eqref{eq:bounded-total-depth-sse-choice}}{>}0$.

\subsection{Regularization}
\label{app:transfer-regular}
Fix a target factor $r$ and a larger factor $r'$ in the same
hardness range. For NP-hardness of total depth, take
$1\le r<r'<63/62$.
Choose an odd $d$ larger than the maximum degree of the
bounded-degree source family. We may increase $d$ further, then
require a sufficiently large source order $n$, using the
enlargements described in Section~\ref{subsec:bounded-degree-consequences}.
Write $G$ for a graph in this family, $n=|V(G)|$, and $R$ for its
$d$-regular output from Proposition~\ref{prop:regularization}.

For treedepth and treewidth, Proposition~\ref{prop:regularization}
changes the promise bounds by at most the fixed constant $d+2$.
Writing $Y$ for their YES upper bound, $Y=\Theta(n)$, so this is $o(Y)$ as the source
order grows. Lemma~\ref{lem:gap-preservation} preserves the
factor-$r$ gap.

A chain on all $n$ vertices is always an elimination forest, so
\begin{equation}\label{eq:total-depth-order-upper}
 \totd(G)\le\frac{n(n+1)}2.
\end{equation}
For total depth, Proposition~\ref{prop:regularization} gives
\begin{equation}\label{eq:regularization-normalized-error}
0\overset{\text{Prop.~\ref{prop:regularization}}}{\le}
 \frac{\totd(R)}{d+2}-\totd(G)
 \overset{\text{Prop.~\ref{prop:regularization}}}{\le}
 \frac{\totd(G)}{d+2}+\frac{d+3}{2}n \overset{\eqref{eq:total-depth-order-upper}}{\le}
 \frac{n(n+1)}{2(d+2)}+\frac{d+3}{2}n.
\end{equation}
For the UGC and SSE total-depth families, let $Y=\Theta(n^2)$ be
the YES upper bound. For any fixed
$\epsilon>0$, choose $d$ large enough that the first term in the
last bound is at most $\epsilon Y/2$ for all sufficiently large $n$,
then choose $n$ large enough that the second term is at most
$\epsilon Y/2$. The normalized error is then at most $\epsilon Y$,
and the NO bound does not decrease. Choosing
$\epsilon<(r'-r)/(r+1)$, Lemma~\ref{lem:gap-preservation}
preserves the factor-$r$ gap.

For the NP total-depth family, let $G=H'$ be the sparse instance
from Corollary~\ref{cor:total-depth-bounded-degree}, and write
$s=|V(G)|=99m+3M$, with $m$ and $M$ as defined
there. In this calculation, $\OPT(J)=\totd(J)$ for every graph $J$.
Appendix~\ref{app:transfer-total-np} gives
$\OPT(G)=\totd(H')
\overset{\eqref{eq:bounded-total-depth-np-sparse}}{\geq}Mm/2=:L$.
Since $M=\Theta(m)$,
we have $s=\Theta(m)$ and $L=\Omega(s^2)$. Choose $\epsilon>0$
such that
\begin{equation}\label{eq:regular-total-depth-np-factor}
 r(1+\epsilon)<r'.
\end{equation}
Use
\eqref{eq:regularization-normalized-error} with $n=s$, and fix
$c>0$ such that
\begin{equation}\label{eq:regular-total-depth-np-scale}
 L\geq cs^2
\end{equation}
throughout this family. For $s\geq1$, the first term is at most
$s^2/(d+2)\overset{\eqref{eq:regular-total-depth-np-scale}}{\leq}
L/(c(d+2))$. Choose an odd $d$, also larger than the maximum degree
of $G$, and then a source order $s$ such that
\begin{equation}\label{eq:regular-total-depth-np-choices}
 \frac1{c(d+2)}\leq\frac\epsilon2,\qquad
 s\geq\frac{d+3}{c\epsilon}.
\end{equation}
For these choices, each term in the last bound
of~\eqref{eq:regularization-normalized-error} is at most
$\epsilon L/2$, using~\eqref{eq:regular-total-depth-np-scale}
for the second. Therefore
\begin{equation}\label{eq:regular-total-depth-np-error}
\frac{\totd(R)}{d+2}-\totd(G)
 \overset{\eqref{eq:regularization-normalized-error},\,
             \eqref{eq:regular-total-depth-np-scale},\,
             \eqref{eq:regular-total-depth-np-choices}}{\leq}\epsilon L,\qquad \OPT(R)\overset{\eqref{eq:regular-total-depth-np-error}}{\leq}
 (d+2)\OPT(G)+(d+2)\epsilon L.
\end{equation}
The projection in Proposition~\ref{prop:regularization} converts
every forest of $R$ of cost $z$ into one of $G$ of cost at most
$z/(d+2)$. In particular, if $z\leq r\OPT(R)$, the projected
forest $\overline F$ satisfies
\[
\totd(\overline F)
 \overset{\text{Prop.~\ref{prop:regularization}}}{\leq}
 z/(d+2)\leq r\OPT(R)/(d+2)
  \overset{\eqref{eq:regular-total-depth-np-error}}{\leq}
 r\OPT(G)+r\epsilon L
 \leq r(1+\epsilon)\OPT(G).
\]
The last inequality uses $L\leq\OPT(G)$. Thus an $r$-approximation
on regular graphs yields factor $r(1+\epsilon)
\overset{\eqref{eq:regular-total-depth-np-factor}}{<}r'$
on the source family.

\subsection{Subdivision}
\label{app:transfer-subdivision}
Let $G$ be a $d$-regular source graph, and let $G_\ell$ be obtained
by replacing each edge of $G$ with a path of length $\ell$.
Use the structural bounds proved in
Corollary~\ref{cor:subdivision-hardness}:
\[
s\totd(G)\le\totd(G_\ell)\le3s\totd(G)+Cn,
 \qquad s=1+(\ell-1)(d-1)/2,\qquad C=h(\ell-1)d/2,
\]
where $h=\lceil\log_2\ell\rceil$ and $n=|V(G)|$.
First fix a target factor
$r$ and any fixed $r'>r$, and choose a regular SSE source family with
NO bound $Z$ and YES upper bound $Y$ satisfying $Z>3r'Y$.
This fixes the odd degree $d$. Fix the required even subdivision
length $\ell$, and hence $h,s,C$, before enlarging the source.
Choose $c_Y>0$ such that $Y\geq c_Yn^2$ throughout this family,
and take $n>rC/[3s(r'-r)c_Y]$. These choices give
\begin{equation}\label{eq:subdivision-sse-choices}
 Z>3r'Y,\qquad 3s(r'-r)Y>rCn.
\end{equation}
The subdivided
instance has YES bound at most $3sY+Cn$ and NO bound at least $sZ$.
Therefore the NO lower bound minus $r$ times the YES upper bound is
at least
\[
sZ-r(3sY+Cn)
 \overset{\eqref{eq:subdivision-sse-choices}}{>}
 3s(r'-r)Y-rCn
 \overset{\eqref{eq:subdivision-sse-choices}}{>}0
\]
Such orders are obtained from
parallel host-edge copies as described in
Section~\ref{subsec:bounded-degree-consequences}.
